\section{Related work}\label{sec:related}

% How each source was read (full text, or metadata only with its content taken from a named
% secondary source) is recorded in docs/research/literature.md. Statements about papers that
% were not read in full are limited to what their titles and abstracts say or are attributed
% to the secondary source.

\paragraph{Elastic-foundation models on surface meshes.} The elastic (Winkler) foundation
replaces the elastic half-space by a layer of independent springs, so that the pressure at a
point depends only on the local compression \cite{Johnson1985}; biomechanics has long used it
with one spring per element of a discretized surface \cite{Blankevoort1991,Fregly2003}. In
multibody dynamics its best-known form is Hippmann's polygonal contact model (PCM)
\cite{Hippmann2004,Hippmann2024}: every triangle of one surface whose centroid lies inside the
other body is a contact element, with a force equal to a layer stiffness in N/m$^3$ times its
area and its penetration along its own normal, an area-related damping and a regularized
friction per element; it has been applied to geometries as demanding as rolling bearings with
defects \cite{Vehvilainen2024}. Simbody's elastic foundation force
\cite{Sherman2011,SimbodyEFF} places one spring at the centre of each face; its friction is a
continuous function of the slip velocity and, as its documentation says, cannot represent
stiction. The normal law of this paper is the same as these, with the same unit of stiffness.
What differs is the quadrature: these models sample the foundation once per face, so whole
faces switch on and off and planar regions must be meshed finely \cite{Hippmann2024}; in
exchange they handle non-convex shapes directly and measure the depth along local normals,
which suits curved conforming contact.

\paragraph{Volumetric models.} Gonthier, McPhee, Lange and Piedb{\oe}uf \cite{Gonthier2005}
integrate the Winkler foundation over a contact surface that is flat or nearly so, and obtain
the normal force, the rolling resistance and the friction in terms of the volume, the centroid
and the inertia tensor of the overlap of the undeformed bodies, with a hysteretic damping factor
(our $h$) from the same authors' regularized point model \cite{Gonthier2004}. Gonthier's thesis
\cite{Gonthier2007thesis} is the complete statement: Coulomb friction with the sliding force and
the spinning torque sharing one budget through an average surface velocity, and a bristle model
for stick--slip \cite{Gonthier2007detc}; the overlap is approximated there with voxels. The
model has been used with closed-form volumes of simple shapes, validated experimentally for the
normal force \cite{Boos2013}, applied to foot--ground contact \cite{Shourijeh2015,Brown2018} and
to wheel--soil interaction \cite{Petersen2013}, and given a damping based on material
properties \cite{Roy2012}. The 2018 thesis behind this paper \cite{Vilela2018} evaluated it on
the exact intersection of triangle meshes and compared it with inner sphere trees
\cite{Weller2009}. Other methods use the overlap without the foundation's derivation: according
to the classification of Elandt et al.~\cite{Elandt2019}, the models of Hasegawa and Fujii
\cite{Hasegawa2003} and of Luo and Nahon \cite{Luo2006,Luo2011} take the direction of the patch
from the vector area of the boundary of the overlap and a force that grows with its volume, and
Wakisaka and Sugihara \cite{Wakisaka2017} combine a volumetric intersection with kinematic
constraints; the same authors observe that all these methods ``rely on a single representative
patch normal, so they do not extend well to non-convex contact''. That sentence applies to the
two laws of this paper too. In computer graphics, layered depth images give the intersection
volume and its gradient \cite{Faure2008,Allard2010}, and Bouchard et al.~\cite{Bouchard2015}
build a constraint with rolling and spinning friction from the moments of the intersection
volume.

\paragraph{Hasegawa and Sato's contact volume model.} The closest prior art to the
\law{distributed} law is the model of Hasegawa and Sato \cite{Hasegawa2004}. They decompose
polyhedra into convex pieces, build the exact intersection of each pair of pieces
\cite{Muller1978}, define one contact plane and normal per pair, and integrate in closed form,
over each triangle of the intersection polyhedron, a distributed spring and damper whose depth
is linear there. They distribute Coulomb friction in the same way but, in their words,
``analytical integration of this equation is difficult'', so they interpolate the dynamic
friction from the three vertices of each triangle; static friction is a distributed spring
integrated into a translational and a rotational spring. The \law{distributed} law of this
paper has the same architecture. It differs in the normalization, which is the Winkler one
(N/m$^3$, so that the force of a flat contact is $K V$, as in Gonthier's model and in
hydroelastic contact; theirs is divided by the contact area), in the direction of the traction
(along each face's own normal), in the stiction (stiffness proportional to the load and one
limit surface for force and torque), and in the friction integral, which is done in closed
form. Elandt's thesis \cite{ElandtThesis2022} describes this design as conceptually
straightforward and names friction as its difficulty, citing Hasegawa and Sato.

\paragraph{Pressure-field (hydroelastic) contact.} Elandt, Drumwright, Sherman and Ruina
\cite{Elandt2019} generalize the elastic foundation: a pressure field is defined in advance
inside each body, the contact surface is where the two fields are equal, and the wrench is the
integral over it of the pressure, a rate term of Hunt--Crossley type and Coulomb shear. The
contact surface is curved and lies inside both bodies, so the model needs no representative
normal and extends to non-convex contact. Masterjohn et al.~\cite{Masterjohn2022} adapt it to
discrete-time solvers by turning each polygon of the contact surface into one compliant point
contact at its centroid, which Drake solves with its transition-aware and convex formulations
\cite{Castro2020,Castro2023,Castro2025}. Drake's documentation presents it as a modern
rendition of the elastic foundation and gives, for small penetrations, the normal force
$\kappa V$ with $\kappa = E/H$, the hydroelastic modulus over the depth of the pressure field,
citing Gonthier \cite{DrakeHydroParams}. With $E = K H$ the model is therefore the same
foundation as ours near the surface. It needs a volume mesh or an analytic field for every
compliant body and has no contact surface between two rigid geometries; in its discrete form
the friction torque within one polygon is not represented and converges with the refinement of
the contact surface \cite{DrakeHydroGuide}.

\paragraph{Point contact.} The continuous point-contact models descend from Hertz
\cite{Hertz1882} with a dissipative term \cite{HuntCrossley1975,Lankarani1990,Flores2011};
Flores et al.~\cite{Flores2023} review contact--impact events with friction. Project Chrono
\cite{Tasora2016} offers a penalty formulation (SMC), with a tangential model whose dependence
on the displacement history has been studied \cite{Fleischmann2016}, and a complementarity
formulation (NSC) in the line of time-stepping schemes for rigid bodies with Coulomb friction
\cite{Stewart1996,Anitescu1997}, solved as a cone complementarity problem
\cite{Anitescu2010,Mazhar2015}. MuJoCo \cite{Todorov2012,Todorov2014} formulates contact as a
convex optimization with soft constraints, with torsional and rolling friction as coefficients
of each contact \cite{MuJoCoDocs}, and Drake has a compliant point contact next to the
hydroelastic one \cite{DrakePointContact}. Yamane and Nakamura \cite{Yamane2006} give a stable
penalty model at the vertices of polygonal objects. In all of these a face-to-face contact is a
small set of points, usually at the boundary of the face.

\paragraph{Conformal contact in joints with clearance.} The multibody literature on revolute
joints with clearance treats the geometry of the \scn{journal} and \scn{bearing} scenarios of
this paper, a pin in a hole of nearly its own radius, with finite-element and analytic
references for the cylindrical contact force \cite{Liu2007,Pereira2011,Pereira2015}. These are
force laws for one geometry, with references that this paper lacks
(Section~\ref{sec:limitations}).

\paragraph{Friction of a patch.} The coupling between the sliding and the spinning friction of
a patch was described by Contensou \cite{Contensou1963}, developed analytically by Zhuravlev
\cite{Zhuravlev1998} and Kireenkov \cite{Kireenkov2008}, and cast as a set-valued force law by
Leine and Glocker \cite{Leine2003}. In manipulation the same physics is the limit surface of a
planar sliding contact \cite{Goyal1991,Mason1986}, with ellipsoidal \cite{Howe1996},
soft-finger \cite{Xydas1999,Ciocarlie2007} and convex polynomial \cite{Zhou2018}
approximations. Jia and Erdmann \cite{Jia1997,Jia1999} give the friction force and moment on a
sliding polygon under uniform pressure in closed form, edge by edge, about the instantaneous
centre of rotation; the integrals are of the family of the potential integrals for uniform and
linear source densities on polygons of boundary-element methods \cite{Wilton1984}. The closed
form of Section~\ref{sec:friction} extends the result of Jia to a linear pressure on a triangle
and uses it as the friction quadrature of a compliant contact law. We did not find that
combination in the contact-simulation literature; the search did not cover the Russian-language
tribology literature thoroughly. Kudra and Awrejcewicz \cite{Kudra2011,Kudra2015} approximate
the integral friction of a patch coupled with rolling resistance and validate it
experimentally, and Marques et al.~\cite{Marques2016,Marques2019} survey friction models for
multibody dynamics, including rolling friction. Regularized friction holds a contact at rest
with a spring that stands for the pre-sliding deflection: the Dahl model \cite{Dahl1976}, LuGre
\cite{CanudasDeWit1995,Olsson1998} and the elasto-plastic model \cite{Dupont2002}; in the tyre
form of LuGre \cite{CanudasDeWit1999} and in Gonthier's bristle \cite{Gonthier2007thesis} the
stiffness is proportional to the load, and the 2018 thesis used the constant-stiffness bristle
of Dopico et al.~\cite{Dopico2011}. The load-proportional stiction of
Section~\ref{sec:stiction} is therefore standard practice and not a contribution of this paper.

\paragraph{Benchmarks and comparisons of simulators.} Gonz\'alez et
al.~\cite{Gonzalez2006,Gonzalez2010} proposed a benchmarking system for multibody simulation
with standardized problems and reference solutions, continued by the library of the IFToMM
technical committee \cite{IFToMMBenchmark} and by recent problems \cite{Ruggiu2023}; these
target mechanisms and formulations, and we did not find conforming frictional contact among
them (the catalogue was not audited problem by problem). Penalty and constraint formulations
have been evaluated against each other \cite{Drumwright2011}, and solvers for the discrete
frictional contact problem compared on a large collection \cite{Acary2018}. In robotics,
engines are compared by their consistency under time-step refinement and by speed--accuracy
curves \cite{Erez2015}, on analytic trajectories and conservation laws \cite{SimBenchmark}, by
predictability \cite{Chung2016}, against measured impacts \cite{Fazeli2017,Acosta2022} or
synthetic data \cite{Guo2023}, by classes of models \cite{Horak2019} and against criteria
derived from the physical laws of contact \cite{LeLidec2024}; Erleben \cite{Erleben2018} gives
a methodology for assessing how contact points are generated from meshes, and Collins et
al.~\cite{Collins2021} review simulators more broadly. Hydroelastic contact was validated by
its authors against Hertz theory and with convergence studies
\cite{Elandt2019,Masterjohn2022}. Two practices of this literature shape
Section~\ref{sec:benchmark}: the reference must not embed the model under test, and accuracy
must be reported against cost and under refinement, not at a single step size.
